\documentclass[a4paper,11pt]{article}
\usepackage{exam-style-842-unit}
\examinehead{842 信息化概论 · 第09单元练习册}{练习一\quad 智慧农业与精准农业}
\begin{document}

\unitcover{842 信息化概论 · 单元练习册}{第09单元\quad 信息化与农业应用}{练习一\quad 智慧农业与精准农业（共 100 分）}{约 70 分钟}{智慧农业的概念、特征与支撑技术；精准农业的概念、实施过程与技术分类；数字农业、精准农业、智慧农业的关系辨析；我国智慧农业发展的现状与主要问题。}

\begin{knowmap}
\textbf{概念递进：}数字农业（农业要素与过程数字化）$\rightarrow$ 精准农业（按空间差异变量投入）$\rightarrow$ 智慧农业（感知—决策—执行闭环、全链条智能化）\\[2pt]
\textbf{智慧农业特征：}感知化 · 数字化 · 网络化 · 智能化 · 精准化 · 全链条化\\[2pt]
\textbf{支撑技术：}物联网（全面感知·可靠传输）· 大数据与云计算（存储·计算）· 人工智能（分析决策）· 3S（定位·遥感·空间分析）· 智能装备（无人机·智能农机）· 区块链（溯源与信任）\\[2pt]
\textbf{精准农业实施链：}信息采集 $\rightarrow$ 差异分析（生成处方图）$\rightarrow$ 决策生成 $\rightarrow$ 变量作业 $\rightarrow$ 效果评价与反馈\\[2pt]
\textbf{精准农业技术分类：}信息获取（RS·GPS·传感器）· 信息处理（GIS·决策支持系统·专家系统）· 变量作业（变量施肥·变量喷药·智能灌溉）· 支撑保障（网络·数据标准·社会化服务）\\[2pt]
\textbf{我国现状：}示范引领、点上突破，整体处于起步与示范推广阶段；短板在基础设施、核心装备、数据标准、人才与推广服务。
\end{knowmap}

\begin{keyterms}
\textbf{智慧农业}：以物联网、大数据、人工智能等为支撑，实现农业生产全过程感知、决策、控制、服务智能化的现代农业形态。\\
\textbf{精准农业}：基于农田空间差异，综合运用 3S、传感器与变量作业技术，实现定位、定时、定量按需投入的农业生产方式。\\
\textbf{数字农业}：把农业要素、生产过程与管理服务数字化的农业发展阶段，是精准农业与智慧农业的基础。\\
\textbf{农业物联网}：通过传感器、RFID 等感知设备与网络，实现农业对象识别、定位、跟踪、监控和管理的网络。\\
\textbf{农业大数据}：农业生产、经营、管理、服务中形成的多源海量数据及其采集、存储、分析与应用。\\
\textbf{变量作业}：依据地块差异图（处方图）实时调整施肥、喷药、灌溉等投入量的作业方式。\\
\textbf{处方图}：在 GIS 中把地块差异与农艺决策落到具体空间位置、用于指导变量作业的图件。\\
\textbf{决策支持系统}：以数据、模型和知识辅助管理者解决半结构化决策问题的人机交互系统。\\
\textbf{智能农机}：具备自动导航、作业监测与变量作业能力的农业机械。\\
\textbf{农业无人机}：用于农田遥感监测、植保喷洒、播种等作业的低空飞行平台。\\
\textbf{全产业链}：涵盖产前投入品、产中种养管理、产后加工流通销售的完整链条。\\
\textbf{农业社会化服务}：由专业服务组织为农户提供农技、农机、植保、烘干等服务的经营形式。
\end{keyterms}

\begin{contrast}{数字农业 / 精准农业 / 智慧农业（每行依次为：核心特征·技术重点·所处阶段）}
\citem{数字农业}{农业要素、过程与管理服务数字化}{数据库、网络、数字采集与管理系统}{农业信息化的数字化基础阶段}
\citem{精准农业}{依据农田空间差异按需投入、变量作业}{3S、传感器、处方图与变量作业装备}{数字农业在生产环节的深化}
\citem{智慧农业}{全链条感知—决策—执行闭环、自主智能}{物联网、大数据、人工智能、智能装备}{农业信息化发展的高级阶段}
\end{contrast}

\newpage
\begin{contrast}{传统农业 / 精准农业（每行依次为：管理单元·投入依据·作业方式）}
\citem{传统农业}{以地块或片区为管理单元}{凭经验统一投入，肥水药“一刀切”}{均一作业、大水大肥}
\citem{精准农业}{以格网或分区形成的管理单元}{依据空间差异与处方图按需投入}{变量施肥、变量喷药、精准灌溉}
\end{contrast}

\begin{contrast}{信息获取 / 信息处理 / 变量作业（每行依次为：环节任务·代表技术·输出结果）}
\citem{信息获取}{采集土壤、作物、环境的空间差异数据}{遥感、GPS／北斗、传感器、无人机}{地块差异数据与专题图层}
\citem{信息处理}{分析差异、生成农艺决策与作业处方}{GIS 空间分析、决策支持系统、专家系统}{管理单元划分与处方图}
\citem{变量作业}{按处方图精确实施农事作业}{变量施肥机、变量喷药机、智能灌溉、自动驾驶农机}{按需投入的精准作业}
\end{contrast}

\qtypename{一、单项选择题}{10 题，每题 3 分，共 30 分}
\q 智慧农业主要以（\quad）为核心驱动，实现农业生产全过程的感知、决策与控制。
\choicesline{A. 劳动力投入}{B. 数据与人工智能}{C. 化肥农药}{D. 土地规模}

\q 精准农业的核心思想是（\quad）
\choices{A. 大面积统一施肥}{B. 根据农田空间差异按需投入}{C. 完全依靠人工经验}{D. 只追求产量最高}

\q 精准农业实施过程的第一步是（\quad）
\choicesline{A. 变量作业}{B. 决策生成}{C. 信息采集}{D. 效果评价}

\q 在精准农业中，把地块差异与施肥建议落到具体空间位置、用于指导变量作业的图件是（\quad）
\choicesline{A. 地形图}{B. 处方图}{C. 气象图}{D. 现状图}

\q 下列技术中，主要用于获取农田空间差异信息的是（\quad）
\choicesline{A. 变量施肥机}{B. 遥感与传感器}{C. 电商平台}{D. 专家系统}

\q 智慧农业中承担“全面感知”任务的技术主要是（\quad）
\choicesline{A. 物联网感知技术}{B. 区块链}{C. 云计算}{D. 虚拟现实}

\q 下列不属于智慧农业支撑技术的是（\quad）
\choicesline{A. 物联网}{B. 大数据}{C. 人工智能}{D. 人畜力耕作}

\q 关于数字农业、精准农业与智慧农业的关系，表述正确的是（\quad）
\choices{A. 三者是同一概念的不同叫法}{B. 智慧农业是数字农业和精准农业发展到高级阶段的形态}
{C. 精准农业在内涵上包含智慧农业}{D. 数字农业与信息技术无关}

\q 我国智慧农业当前总体上处于（\quad）
\choices{A. 全面普及阶段}{B. 起步与示范推广阶段}{C. 完全成熟阶段}{D. 尚未起步阶段}

\q 精准农业能够减少化肥农药用量的根本原因是（\quad）
\choices{A. 减少种植面积}{B. 按地块差异精准投入、避免过量}{C. 降低产量目标}{D. 用人工替代机械}

\qtypename{二、多项选择题}{4 题，每题 4 分，共 16 分；每题有两个或两个以上正确项}
\q 智慧农业的基本特征包括（\quad）
\choicesline{A. 感知化}{B. 数字化}{C. 智能化}{D. 粗放化}

\q 精准农业的主要支撑技术包括（\quad）
\choicesline{A. 遥感与传感器}{B. 地理信息系统}{C. 全球定位系统}{D. 人工经验估算}

\q 精准农业的实施过程包括（\quad）
\choicesline{A. 信息采集}{B. 差异分析}{C. 决策生成}{D. 变量作业}

\q 我国发展智慧农业面临的主要问题有（\quad）
\choices{A. 农村信息基础设施仍较薄弱}{B. 高端传感器与智能装备依赖进口}
{C. 数据标准与共享机制不健全}{D. 农业劳动力严重过剩}

\qtypename{三、判断题}{6 题，每题 2 分，共 12 分；正确的画“√”，错误的画“×”并写出改正后的表述}
\q 智慧农业就是在农田安装摄像头、用手机远程查看画面。\quad（\quad）\quad 改正：\fixblank

\q 精准农业的核心是根据农田空间差异实施变量投入。\quad（\quad）\quad 改正：\fixblank

\q 精准农业实施时，信息采集后可直接进行变量作业，不需要差异分析和决策生成。\quad（\quad）\quad 改正：\fixblank

\q 数字农业、精准农业和智慧农业是同一概念的不同叫法。\quad（\quad）\quad 改正：\fixblank

\q 智慧农业是数字农业和精准农业发展到高级阶段的产物。\quad（\quad）\quad 改正：\fixblank

\q 变量作业必须以地块差异信息和处方图为依据。\quad（\quad）\quad 改正：\fixblank

\qtypename{四、填空题}{6 空，每空 1 分，共 6 分}
\q 智慧农业的技术特征通常概括为感知化、数字化、网络化、\fillblanks[2.0cm]和精准化。

\q 精准农业的实施过程为：信息采集 $\rightarrow$ \fillblanks[2.2cm] $\rightarrow$ 决策生成 $\rightarrow$ 变量作业。

\q 精准农业的核心思想是“因地制宜、\fillblanks[2.2cm]”。

\q 在精准农业技术体系中，变量施肥机、变量喷药机属于\fillblanks[2.0cm]技术。

\q 我国智慧农业整体上处于起步和\fillblanks[2.0cm]阶段，尚未全面普及。

\q 智慧农业以\fillblanks[1.6cm]为关键生产要素，以人工智能为核心驱动。

\newpage
\qtypename{五、简答题}{4 题，每题 6 分，共 24 分}
\q 什么是智慧农业？简述智慧农业的基本特征与主要支撑技术。\anslines{2.4cm}

\q 什么是精准农业？简述精准农业的实施过程（四个环节）。\anslines{2.4cm}

\q 简述精准农业技术的主要分类，并各举一种代表技术。\anslines{2.4cm}

\q 简述我国智慧农业发展的现状与存在的主要问题。\anslines{2.4cm}

\newpage
\qtypename{六、综合题}{1 题，共 12 分}
\q 某县以水稻和柑橘为主导产业，计划建设智慧农业示范区：田间布设土壤墒情与气象传感器，无人机定期巡田，建设县域农业大数据平台，配置自动驾驶拖拉机与变量施肥机，并开发手机端应用向农户推送农事建议。请回答：（1）该示范区用到了智慧农业的哪些支撑技术？各承担什么任务？（2）结合精准农业实施过程，说明该示范区如何做到按需施肥、按需灌溉；（3）分析该模式在以小农户为主的地区推广可能遇到的困难。\anslines{6cm}

\end{document}
